\chap{绪论}

\sect{研究背景}

温度历程决定金属热处理中的相变驱动力、组织均匀性和热应力水平。炉温可以直接测量和调节，工件内部温度场通常只能通过少量热电偶间接获得。厚壁工件升温时，表面与心部会形成明显温差。控制系统若只依据炉温或表面温度，容易错判心部到温时刻。因此，能够随控制输入快速推演完整温度场的模型具有直接工程价值。

有限差分、有限体积和有限元方法已经构成热过程分析的标准工具。它们能够处理温变物性和复杂边界，也适合生成高可信度参考解。在线控制提出了另一类计算需求。控制器需要在每个决策时刻比较多条候选炉温序列，并根据新观测重新计算。二维、三维模型中的重复求解成本很高。连续钢坯加热炉的模型化轨迹规划和非线性预测控制研究已经证明，工件温度预测可以直接进入炉温设定与能耗优化\cite{steinboeck2011trajectory,steinboeck2013nmpc}。学习型代理模型可将主要计算转移到离线阶段，为实时推演提供条件。

现有温度场机器学习研究多采用两种任务设置。第一类直接近似给定偏微分方程的解 $T(x,t)$。物理信息神经网络通过自动微分把控制方程和边界条件写入训练目标，在正问题和逆问题中形成了统一框架\cite{raissi2019pinn,karniadakis2021piml}。第二类从有限测点恢复当前或历史温度场。相关方法已用于钢材加热、水淬全场重构和虚拟热传感器\cite{sun2025steel,zhao2022quenching,go2023sensor}。这些任务推动了温度场学习，但工艺决策还需要一个动作条件化的状态转移模型。模型应接收当前温度场和下一段炉温指令，并能递归计算控制改变后的未来状态。

世界模型通常指能够依据当前状态和动作预测未来环境状态的内部动力学模型，并可供策略学习或规划调用\cite{ha2018worldmodels}。本文采用面向热过程的确定性定义，将世界模型写为
\begin{equation}
  \bm{s}_{n+1}=F_{\theta}(\bm{s}_n,a_n,\bm{p}),
  \label{eq:wm_definition}
\end{equation}
其中 $\bm{s}_n$ 是离散温度场，$a_n$ 是炉温动作，$\bm{p}$ 包含几何、物性和边界参数。式\eqref{eq:wm_definition}可以连续调用，所以同一模型能够评价多条候选工艺轨迹。本文直接使用可观测温度场作为状态，不引入图像生成任务中常见的潜变量和概率解码器。“世界模型”在这里强调动作条件化的递归推演及其规划用途。就数学形式而言，它属于学习型非线性状态空间模型。若模型只执行一次前向温度预测，它可归入代理动力学模型；本文的命名来自它在状态估计和控制器内部承担的环境演化与反事实动作评价功能，不表示提出了新的通用神经网络架构。

\sect{相关研究}

\subsect{物理信息学习与温度场预测}

PINN 将方程残差、初边值条件和观测误差共同用于网络训练\cite{raissi2019pinn}。它减少了对标注数据的依赖，并可在同一计算图中处理未知参数。本文的训练主体是数值轨迹监督下的神经动力学模型，离散方程残差用于正则化，因此不属于主要依靠配点残差求解偏微分方程的经典 PINN。训练效果会受到损失尺度、梯度冲突和采样分布影响\cite{wang2021gradient,krishnapriyan2021failure}。物理残差较低并不自动对应更小的场误差。温度场应用还要面对长时递归误差、控制分布变化和边界参数漂移。评价口径需要覆盖完整轨迹。

Sun 等将物理信息网络用于热轧钢加热温度分布预测\cite{sun2025steel}。Zhao 等利用稀疏测温信息恢复水淬过程的时空温度场，并估计热边界\cite{zhao2022quenching}。Go 等研究了有限温度观测下的虚拟热传感器和未知热流估计\cite{go2023sensor}。Tian 等进一步使用物理信息递归网络预测金属增材制造中的长时温度场，重点处理递归误差和工艺变化\cite{tian2025longhorizon}。这些研究证明了物理信息学习在金属热过程中的可行性。本文关注炉温动作驱动的一维热处理轨迹，并把动态边界可辨识性纳入状态转移和控制链路。

\subsect{神经算子与学习型动力学模型}

DeepONet 和 Fourier Neural Operator 学习输入函数或参数到解场之间的映射\cite{lu2021deeponet,li2021fno}。物理信息 DeepONet 进一步把控制方程用于算子训练，可面向一类初边值条件快速求解\cite{wang2021pideeponet}。这类模型适合跨网格、跨载荷或跨几何问题。本文当前对象为固定的一维 41 节点网格。轻量残差多层感知机已经能够覆盖状态维数，并具有较低推理开销。研究重点因此放在动力学训练、物理约束和边界可辨识性上。

学习型物理模拟器常采用自回归推演。一步监督训练只接触真实状态，部署时却持续接收模型生成状态，两种输入分布的偏移会放大长期误差。List 等的研究表明，时间展开方式会影响瞬态动力学的轨迹精度\cite{list2025unrolled}。Sholokhov 等把物理配点约束写入潜空间神经常微分方程，用于高维动力学的降阶预测与控制\cite{sholokhov2023pinode}。本文比较一步与五步展开，并以完整验证轨迹选择检查点。五步设置在 ID 测试中取得较低误差，控制曲线 OOD 则呈现精度权衡。因此，展开训练在本文中作为与部署形式一致的训练设置，不单独承担 OOD 优势主张。

传统降阶模型也可形成面向控制的热动力学表示。Kloeser 等将热成形有限元模型降至低维状态，并结合扩展卡尔曼滤波从稀疏测点估计时空温度和参数\cite{kloeser2021rom}。该路线依赖显式降阶基和降阶方程。本文直接学习固定网格上的非线性状态转移，并通过离散能量残差约束轨迹。两类方法都可嵌入模型预测控制\cite{rawlings2017mpc}，共同结构是快速动力学模型与在线状态估计。本文进一步处理逐时动态对流辐射边界的输入表示，并在相同动作搜索配置下比较神经动力学与 BDF 动力学。

\subsect{动态热边界辨识}

高温炉内钢材加热同时受对流和辐射影响。辐射贡献随表面温度快速变化，钢材氧化状态还会改变发射率。Kim 的加热炉模型说明了辐射边界和钢坯内部瞬态导热的耦合关系\cite{kim2007furnace}。边界参数往往无法直接测量，逆热传导方法需要从表面温度或热流响应中估计载荷。Frankel 和 Arimilli 研究了瞬态表面测温下对流与辐射载荷的反演\cite{frankel2007loads}。钢坯多表面热流反演、序贯正则化和神经正向模型也已用于在线边界估计\cite{li2015slabflux,chen2019online,mohebbi2021htc}。

对流系数 $h(t)$ 和发射率 $\varepsilon(t)$ 若同时按时间自由变化，单个时间步的温度响应无法提供两个独立观测方向。将二者直接作为网络输入，会把等价热边界表示成多组不同参数。本文对离散边界方程进行灵敏度分析，确定可观测组合 $\Heff(t)$，并据此设计世界模型输入和在线状态估计器。

\subsect{工业炉模型与热处理数字孪生}

钢铁生产中的在线温度模型通常连接二级控制系统。Steinboeck 等以离散非线性热模型规划钢坯热流和炉温轨迹，并在连续加热炉上实现非线性模型预测控制\cite{steinboeck2011trajectory,steinboeck2013nmpc}。Yi 等构建炉体与钢坯双模型跟踪系统，持续计算炉内钢坯温度场\cite{yi2017tracking}。Kavak 和 Yalcin 使用实际运行数据辨识 12 区步进梁式加热炉，并以扰动观测器表示未建模热损失\cite{kavak2023furnace}。这些研究说明，工件内部温度、动态热损失和上层炉温设定是工业部署中的核心变量。

数字孪生强调物理对象、数字模型、在线数据和决策服务之间的持续连接\cite{tao2018digitaltwin,stark2019digitaltwin}。在材料加工领域，Derouiche 等通过降阶与回归模型预测感应淬火中的温度和奥氏体演化\cite{derouiche2021induction}。Gong 等构建端淬温度--组织耦合降阶模型和在线仿真平台\cite{gong2026heatrom}。Midaoui 等利用有限元快照、POD 和红外表面测温重建 Jominy 端淬试样的三维温度场，并进行了数值与实验验证\cite{midaoui2026jominy}。本文与这些工作的共同目标是获得可在线更新的热状态。本文进一步把炉温作为显式动作输入，并将动态边界的可辨识表示、后验估计和反事实动作推演接入同一递归模型。

\sect{研究内容与贡献}

本文围绕受控热处理过程建立一条完整方法链。主要贡献如下。

\begin{enumerate}
  \item 构建面向金属热处理温度场的受控递归世界模型。模型显式接收炉温动作和热参数，采用完整轨迹 RMSE、尾部误差和物理一致性评价递归推演质量。
  \item 建立覆盖控制曲线、换热系数、发射率和温变物性的 OOD 实验矩阵。结果给出物理约束的适用条件。其主要收益集中在组合参数外推、冷却轨迹、尾部误差和离散能量一致性。
  \item 分析动态对流辐射边界在单步状态转移下的局部可辨识性，并以 $\Heff$ 重新参数化世界模型输入。基于该表示构建因果观测器、集合卡尔曼滤波和滚动控制。交叉求解器、稀疏测点和风险感知控制实验验证了方法链的有效性。
\end{enumerate}

全文共六章。第二章给出物理模型和数值基准。第三章介绍世界模型、物理损失和边界参数化。第四章分析展开长度、控制曲线 OOD 和参数 OOD。第五章研究动态边界辨识、跨求解器验证、部分可观测状态估计和闭环控制。第六章总结研究结论并给出后续研究方向。
